% Retrospective synthesis; numerical sources and detailed chronology are kept in
% research/progress-review-2026-11-06/{experiment-history,data-method-audit}.md.
% Historical JSON: G2, learning-curve P0/P1, R0/R1/R2/R3/R6, scan-sat-merge, P3.
\section{Research Progress and Data Evolution}
\label{sec:research-progress}
The investigation has changed its experimental focus as evidence accumulated. Coverage and
stratification address RQ1; representation and head controls address RQ2; source transfer addresses
RQ3; matched adaptation addresses RQ4; and grouped evaluation, errors, calibration and transfer
boundaries address RQ5. This section
reconstructs that reasoning from recorded protocols and results; it does not retrospectively assign
preregistered status to exploratory questions. Historical development experiments and the current
full-PSI campaign remain distinct. The main milestones are summarised in
Table~\ref{tab:research-progress-history}; Chapter~\ref{chap:evaluation} reports their numerical results.

\begin{table}[htbp]
\centering
\caption{Research progression: the question resolved and the decision it motivated}
\label{tab:research-progress-history}
\small
\begin{tabularx}{\textwidth}{|P{2.0cm}|X|X|}
\hline
\textbf{Stage} & \textbf{Question and recorded finding} & \textbf{Next decision and boundary} \\ \hline
Data audit, 12--15 Oct. & Reconstructed bin-based PSI and distinguished scanning from saturation mutation labels. & Use direct mutation $\Delta$PSI labels and assay-specific WT references; peptide activity is a separate endpoint. \\ \hline
G2 and extensions, 18--19 Oct. & Ridge was a strong saturation baseline; additional nonlinear candidates did not establish an architectural advantage. & Audit training exposure and label availability before adding capacity. The overlapping public checkpoint is unsuitable as clean generalisation evidence. \\ \hline
P0/P1, 22 Oct. & Adding WT-PSI $(3,4]$ training peptides improved prediction on the original low-PSI test population. & Expand training coverage, while testing whether improvement persists outside the original population. Label scarcity remains an explanation to test. \\ \hline
P3 and scan transfer, 22--23 Oct. & Transfer depended on the source target and learner; peptide-score pretraining did not reliably improve mutation prediction. & Keep source/replicate and training recipe explicit; do not pool endpoints merely because both concern degradation. \\ \hline
R0--R3, 22 Oct. & Reproduction and inference parity passed; cleaned Degformer and frozen ESM--Ridge did not establish superiority; the domain gate required assay-scale mapping. & Preserve a reproducible baseline, add matched head controls, and distinguish domain confidence from effect accuracy. \\ \hline
R6, 23 Oct. & Ordered W5 context improved Ridge; the primary ESM--mCNN contrast was inconclusive. A secondary one-hot mCNN scanning-pretraining contrast was positive. & Separate primary and secondary claims; test representation, head and training range on common full-range folds. \\ \hline
Full-PSI v3, 31 Oct.--6 Nov. & Wide evaluation, pointwise/CNN controls, nine frozen pLM recipes and matched frozen--LoRA evaluation were completed where inputs were available. & Interpret development evidence by stratum and metric; final serving export, external validation and structure inputs remain unfinished. \\ \hline
\end{tabularx}
\end{table}

\subsection{From Peptide Activity to Direct Mutation Labels}
The initial data audits established why peptide-level degron scores cannot substitute for mutation
effects. A peptide may receive one activity score, whereas \mbox{saturation} mutagenesis measures many
substitutions at each position. Scanning provides more peptide contexts but only a restricted set of
replacement amino acids. Its substitutions are drawn from R, A and S, rather than all 19 alternatives.
Consequently, sequence context can be informative in scanning cross-validation while transfer to
the saturation substitution distribution remains difficult. This motivated direct training on
saturation labels and separate reporting of the two experimental populations.

The WT reference is part of the endpoint definition. Saturation labels subtract the median of
valid synonymous observations in each half before averaging the two half-specific effects.
Scanning subtracts the measured WT construct of the same replicate. The separate BAG6-WT
condition cannot replace either reference. Nor is the approximately 0.9-PSI offset observed between
the peptidome screen and intrinsic WT measurements an offset between scanning and saturation
mutation labels. These distinctions prevent a change of measurement scale from being mistaken
for biological transfer failure.

\subsection{Why the Study Expanded Training Coverage}
G2 evaluated 91 peptides and 47,847 cells at WT PSI at most 3. The P0 experiment retained those
outer-test cells and added up to 85 peptides, or 44,490 cells, from WT PSI in $(3,4]$ to permitted
training groups. Ridge macro AUROC increased from 0.8082 to 0.8208, with a paired group-bootstrap
95\% interval for the difference of $[0.0061,0.0190]$. This supports that particular expansion on the
original test population. It does not establish performance on the added peptides as test examples.

The associated P1 learning-curve hypothesis was model-dependent. For Ridge, the 100\% versus
50\% contrast had a 95\% interval containing zero and the curve was not monotone; the corresponding
neural learner did support a label-quantity effect. Therefore, the earlier interpretation that label
scarcity alone explained unsuccessful architectures is too strong. Training distribution, noise,
regularisation and representation remain plausible contributors. The full-PSI design subsequently
changed both validation and evaluation coverage to measure performance across the full range.

The wide dataset retains 235 technically valid saturation peptides and 122,412 cells. Its 221 groups
result from the full-proteome mapping and additional homology/sequence unions. Historical counts
of 227 legacy groups or 222 full-proteome groups preceded this union procedure; they do not
describe extra mutation observations. Likewise, the full-range folds are not the original G2 folds,
so an AUROC change between the two populations cannot isolate architectural improvement.

\subsection{Why Matched Representation and Head Controls Were Needed}
R2 added PCA-reduced frozen ESM2-650M features to Ridge and obtained a small, inconclusive
gain. R6 then compared sequence representations and convolutional heads. Its primary
ESM2-35M mCNN--Ridge difference was $+0.0097$, with a 95\% interval of
$[-0.0025,0.0233]$; it did not meet its confirmation criterion. Ordered W5 context did improve
Ridge within its adjusted comparison family. Separately, scanning pretraining improved the
one-hot mCNN over its scratch control by $0.0147$ AUROC, with the adjusted interval
$[0.0027,0.0276]$. That secondary transfer result neither rescues the primary contrast nor
establishes an advantage over W5 Ridge.

The ESM mCNN--single-window CNN difference was only $0.0041$, with a 95\% interval of
$[-0.0005,0.0087]$. The evidence thus motivated a matched pointwise head and broader
evaluation rather than a claim that multiple convolutional scales were necessary. Historical Ridge
and neural fits also used different final refit policies; the current campaign explicitly refits outer
models after inner selection. Recipe, runtime and final fitting policy must accompany any comparison.

\subsection{What the Current Progress Does and Does Not Establish}
The current campaign addresses the new full-range questions using fixed grouped folds, inner
selection and explicit training arms. It provides development evidence on frozen representations,
convolutional heads and a matched LoRA adaptation. It contains no controlled full-fine-tuning arm.
The completed transfer evaluation concerns W5-D alone on a restricted scanning population from
the same source laboratory. Calibration, predictive ranking, regression accuracy and biological
utility remain different outcomes. Neither completion of the runnable graph nor favourable OOF
metrics constitutes final model selection, independent external validation, or a demonstrated
reduction in wet-laboratory cost.
